% Ferejohn 與 Pasquino 第 3 節（28 至 32 頁）：三種體制中憲法審查權之所在。
% 欄序依兩人所說的單一向度，左端議會主權、右端權力分立（28 頁）；Kelsen 型居中係依 30 頁「departed from the
% hierarchical parliamentary model」之描述排列。垂直、水平二詞出自 29 頁「separated horizontally rather than vertically」。
% 色調：審查權所在之機關 accent（terracotta），其餘機關 slate；卡片淡底加左側色條（skill diagram-design「動手前」第 2 條）。
% 各格文字逐一回刊本核過頁次（2026-11-09）。
\definecolor{figRev}{HTML}{B08060}\definecolor{figRevbg}{HTML}{F5ECEB}
\definecolor{figLine}{HTML}{E8D3D1}
% 卡片：\placecard{色}{左}{上}{右}{下}{內容}
\providecommand{\placecard}[6]{%
  \fill[#1bg] (#2,#3) rectangle (#4,#5);
  \fill[#1] (#2,#3) rectangle ($(#2,#5)+(0.6,0)$);
  \node[t] at ($(#2,#3)+(2.2,-1.4)$) {#6};}
\begin{tikzpicture}[qa,x=1mm,y=1mm,
    t/.style={anchor=north west,inner sep=0pt,align=left,execute at begin node={\raggedright\setstretch{1.2}}},
    cell/.style={t,font=\scriptsize,text width=40mm},
    rowlbl/.style={t,font=\scriptsize,text=dimink,text width=17mm},
    link/.style={figGray,line width=0.5pt},
    arr2/.style={figGray,line width=0.6pt,-{Stealth[length=4pt,width=3.2pt]}}]
  % 單一向度
  \draw[figGray,line width=0.5pt,{Stealth[length=3.5pt,width=3pt]}-{Stealth[length=3.5pt,width=3pt]}] (19,6) -- (146,6);
  \node[font=\scriptsize,text=dimink,anchor=south] at (82.5,6.6) {制度位置的單一向度（28 頁）};
  % 欄頭
  \foreach \x/\h/\s in {
      19/議會主權體制/{垂直：行政與法院\\從屬於立法機關（28 頁）},
      62.5/{Kelsen 型模式}/{憲法法院置於司法與\\其他部門之外（30 頁）},
      106/權力分立體制/{水平：三權分立（29 頁）}}{
    \node[t,font=\footnotesize] at (\x,2.5) {\h};
    \node[t,font=\scriptsize,text=dimink] at (\x,-1.8) {\s};
  }
  % 列名與列間淡線
  \node[rowlbl] at (0,-13) {機關之配置};
  \node[rowlbl] at (0,-50) {審查之時點\\與方式};
  \node[rowlbl] at (0,-61.5) {違憲之效果};
  \node[rowlbl] at (0,-73) {例};
  \foreach \y in {-10,-47.5,-59,-70.5}{\draw[figLine,line width=0.4pt] (0,\y) -- (146,\y);}
  % 議會主權：審查在立法機關內部，行政與法院居其下
  \begin{scope}[shift={(19,0)}]
    \placecard{figGray}{0}{-13}{40}{-34}{立法機關}
    \placecard{figRev}{3.5}{-19}{38}{-32}{憲法審查在其內部\\{\scriptsize 上議院、司法委員會、}\\{\scriptsize 法國憲法委員會（28 頁）}}
    \placecard{figGray}{0}{-38.5}{18}{-44.5}{行政}
    \placecard{figGray}{22}{-38.5}{40}{-44.5}{法院}
    \draw[link] (9,-34) -- (9,-38.5);
    \draw[link] (31,-34) -- (31,-38.5);
  \end{scope}
  % Kelsen 型：憲法法院居各部門之外、之上；普通法院停止訴訟並移送
  \begin{scope}[shift={(62.5,0)}]
    \placecard{figRev}{0}{-13}{40}{-29}{憲法法院\\{\scriptsize 審查權集中於此，}\\{\scriptsize 居各機關之上（31 頁）}}
    \placecard{figGray}{0}{-38.5}{10}{-44.5}{立法}
    \placecard{figGray}{12}{-38.5}{22}{-44.5}{行政}
    \placecard{figGray}{24}{-38.5}{40}{-44.5}{普通法院}
    \draw[arr2] (32,-38.5) -- (32,-29);
    \node[t,font=\scriptsize,text=dimink,anchor=north east,align=right] at (30.5,-30.6) {疑法律違憲者\\停止訴訟並移送};
  \end{scope}
  % 權力分立：三權並列，審查權散布於各級法院
  \begin{scope}[shift={(106,0)}]
    \placecard{figGray}{0}{-13}{10}{-19}{立法}
    \placecard{figGray}{12}{-13}{22}{-19}{行政}
    \placecard{figGray}{24}{-13}{40}{-19}{司法}
    \placecard{figRev}{24}{-25.5}{40}{-31.5}{最高法院}
    \placecard{figRev}{24}{-38.5}{40}{-44.5}{下級法院}
    \draw[link] (32,-19) -- (32,-25.5);
    \draw[link] (32,-31.5) -- (32,-38.5);
    \node[t,font=\scriptsize,text=dimink,anchor=north east,align=right] at (22.5,-27.6) {任何法院皆得審查，\\最高法院居其首};
  \end{scope}
  % 各列內容
  \foreach \x/\a/\b/\c in {
      19/{事前、抽象：比較法律\\與憲法之文本（28 頁）}/{法國：提案不得公布；\\公布後不受審查（32 頁）}/{英國、芬蘭、瑞典；法國\\第三至第五共和（28 頁）},
      62.5/{抽象審查之外，兼有\\事後之具體審查（31 頁）}/{廢止法律並重建法律狀態，\\兼為消極與積極立法（30 頁）}/{奧地利、德國、義大利、\\西班牙、葡萄牙（30 頁）；\\東歐（31 頁）},
      106/{於具體案件中為之，\\非抽象、非事前（29 頁）}/{法律不廢止，僅於該案\\拒絕適用（29 頁）}/{美國、挪威（29 頁，註 11）}}{
    \node[cell] at (\x,-50) {\a};
    \node[cell] at (\x,-61.5) {\b};
    \node[cell] at (\x,-73) {\c};
  }
\end{tikzpicture}
